%------------------------------------------------------------------------------------
%  TABLEAU DES HYPOTHÈSES DIAGNOSTIQUES (BOOKTABS & TABULARX PROFESSIONNEL)
%------------------------------------------------------------------------------------
\begin{table}[htbp]
  \centering
  \caption{Synthèse des hypothèses diagnostiques gynécologiques}
  \label{tab:hypothesesDiagnostiques}
  \small
  \begin{tabularx}{\textwidth}{>{\bfseries}p{3cm} X X}
    \toprule
    Diagnostic & Définition & Symptômes \\
    \midrule
    Kyste hémorragique du corps jaune &
    Kyste fonctionnel dépendant du cycle menstruel. &
    Dysménorrhée avec douleur unilatérale lors de la phase lutéale. \\
    \addlinespace
    Endométriome &
    Endomètre ectopique au niveau ovarien. &
    Dysménorrhée, douleurs pelviennes chroniques, dyspareunie profonde, troubles de la fertilité. \\
    \addlinespace
    Kyste dermoïde &
    Tératome, tumeur bénigne d'origine embryonnaire, fréquente entre $16$ et $55$ ans. &
    Asymptomatique ou légère douleur pelvienne, métrorragie, dyspareunie et compression éventuelle. \\
    \addlinespace
    Kyste para-ovarien &
    Kyste bénin avec pour origine le mésosalpinx. &
    Asymptomatique ou symptômes liés à une compression mécanique. \\
    \addlinespace
    Tumeur ovarienne &
    Cancer ovarien malin. Facteurs favorisants : $G_0 P_0$, ménarche précoce, cycles spontanés. &
    \emph{Stade précoce :} paucisymptomatique.\newline
    \emph{Stade avancé :} AEG, perte pondérale, asthénie. \\
    \bottomrule
  \end{tabularx}
\end{table}

\begin{table}[htbp]
  \centering
  \caption{Signes cliniques et remarques associées}
  \label{tab:signesCliniques}
  \small
  \begin{tabularx}{\textwidth}{>{\bfseries}p{3cm} X X}
    \toprule
    Diagnostic & Signes cliniques & Remarques \\
    \midrule
    Kyste hémorragique du corps jaune &
    \begin{itemize}[nosep, leftmargin=*]
      \item Palpation abdominale : sensible ;
      \item Toucher vaginal : masse annexielle palpée, mobilisation utérine douloureuse.
    \end{itemize} &
    Rupture du kyste $\rightarrow$ défense abdominale, signe du cri du Douglas. \\
    \addlinespace
    Endométriome &
    \begin{itemize}[nosep, leftmargin=*]
      \item Palpation abdominale : sensible ;
      \item Toucher vaginal : masse annexielle palpable, adhérences suspectées.
    \end{itemize} &
    Symptômes associés possibles (nycturie, troubles du transit). \\
    \addlinespace
    Kyste dermoïde &
    \begin{itemize}[nosep, leftmargin=*]
      \item Palpation abdominale : normale ;
      \item Toucher vaginal : masse annexielle potentiellement palpable.
    \end{itemize} &
    Torsion ovarienne = urgence chirurgicale $\rightarrow$ douleur aiguë et brutale. \\
    \addlinespace
    Kyste para-ovarien &
    \begin{itemize}[nosep, leftmargin=*]
      \item Palpation abdominale : normale ;
      \item Toucher vaginal : masse annexielle parfois palpable.
    \end{itemize} &
    Évolution le plus souvent asymptomatique. \\
    \addlinespace
    Tumeur ovarienne &
    \begin{itemize}[nosep, leftmargin=*]
      \item Palpation abdominale : signes de carcinose ;
      \item Toucher vaginal : masse annexielle dure, fixée.
    \end{itemize} &
    Diagnostic d'exclusion prioritaire en raison du pronostic vital. \\
    \bottomrule
  \end{tabularx}
\end{table}

